% 3.3 =============== Conditional Expectation and Discrete AR(1) ===============

\section{离散化估计}    \label{sec-discrete-approximation}

我们已经初步了解了马尔可夫性质对于条件期望结构的关键作用——使未来的分布只依赖于当前。
截至目前，我们基本都是在离散空间上研究马尔可夫过程和贝尔曼方程，但我们曾指出，多数情形下外生冲击服从连续过程，
即我们求解的是连续动态规划问题。对此我们也提出了两种总体思路：离散化估计和平滑估计。但无论选择何种总体思路，
\textbf{具体的难点都是处理连续的AR(1)过程，或者说计算积分}。
本节，我们将在对离散化估计进行学习的过程中，掌握马尔可夫近似AR(1)过程和数值积分两种具体计算方法。


% 3.3.1 =====  =====
\subsection{离散化AR(1):Tau-Chen Method}

一个典型的AR(1)冲击形式如下：
\begin{equation}    \label{eq-z-ar1}
    z_{t+1}=(1-\rho) \mu+\rho z_t+\sigma \varepsilon_{t+1}, \quad \varepsilon_{t+1} \sim \operatorname{IID} N(0,1)
\end{equation}
如果选择离散化估计，并且想要在积分计算中直接地利用马尔可夫性质，就必须要求离散化后的外生过程服从马尔可夫链，
并且需要找到对应的转移概率矩阵。因此，对于外生过程的离散化并不像离散化资本那样简单。


离散化后求和仍有V，可以通过离散化后VFI，
也可作为后面函数逼近和Projection的引入，即用函数逼近V


% 3.3.2 =====  =====
\subsection{Solve the Model Numerically}